\chapter{Descripción general}

\section{Perspectiva del producto}
\campo{Inserte aquí el texto}
\orientacion{Indique si el producto es independiente o forma parte de un sistema mayor. Si forma parte de otro sistema, un diagrama que sitúe el producto e identifique sus conexiones facilita la comprensión.}

\section{Funcionalidad del producto}
\campo{Inserte aquí el texto}
\orientacion{Resuma las funcionalidades principales del producto sin entrar en detalles. En ocasiones, esta información puede tomarse de una especificación del sistema de mayor nivel. Organice las funcionalidades de forma que el cliente y los demás interlocutores puedan entenderlas con claridad; puede utilizar descripciones textuales o gráficas.}

\section{Características de los usuarios}
\begin{longtable}{p{3.4cm}p{10.2cm}}
\toprule
\textbf{Característica} & \textbf{Descripción} \\
\midrule
Tipo de usuario & \campo{Inserte aquí el texto} \\
Formación & \campo{Inserte aquí el texto} \\
Habilidades & \campo{Inserte aquí el texto} \\
Actividades & \campo{Inserte aquí el texto} \\
\bottomrule
\end{longtable}
\orientacion{Describa a los usuarios del producto, incluidos su nivel educacional, experiencia y conocimientos técnicos.}

\section{Restricciones}
\campo{Inserte aquí el texto}
\orientacion{Describa las limitaciones que deben considerarse al diseñar y desarrollar el sistema, como metodologías de desarrollo, lenguajes de programación, normas, restricciones de hardware o sistemas operativos.}

\section{Suposiciones y dependencias}
\campo{Inserte aquí el texto}
\orientacion{Describa los factores que, si cambian, pueden afectar a los requisitos. Por ejemplo, puede suponerse que cierto sistema operativo estará disponible para el hardware requerido. Si no lo estuviera, habría que modificar esta especificación.}

\section{Evolución previsible del sistema}
\campo{Inserte aquí el texto}
\orientacion{Identifique mejoras futuras del sistema que puedan analizarse e implementarse posteriormente.}
